\subsection{测试小节三}

本节测试 tikz-cd 交换图.

\subsubsection*{范畴论短交换图}

张量积的泛性质 (图~\ref{cd:tensor}):

\begin{figure}[htbp]
    \centering
    \begin{tikzcd}
        M \times N \arrow[r, "B"] \arrow[d, "\otimes"'] & P \\
        M \otimes N \arrow[ru, "\tilde{B}"', dashed]
    \end{tikzcd}
    \caption{张量积的泛性质}
    \label{cd:tensor}
\end{figure}

自由对象的泛性质 (图~\ref{cd:free}):

\begin{figure}[htbp]
    \centering
    \begin{tikzcd}
        X \arrow[r, "f"] \arrow[d, "i"'] & A \\
        F(X) \arrow[ru, "\tilde{f}"', dashed]
    \end{tikzcd}
    \caption{自由对象的泛性质}
    \label{cd:free}
\end{figure}

自然变换的交换方图 (图~\ref{cd:nat}):

\begin{figure}[htbp]
    \centering
    \begin{tikzcd}
        A \arrow[r, "f"] \arrow[d, "\alpha_A"'] & B \arrow[d, "\alpha_B"] \\
        A' \arrow[r, "f'"] & B'
    \end{tikzcd}
    \caption{自然变换的交换方图}
    \label{cd:nat}
\end{figure}

\subsubsection*{同调代数长交换图}

五引理 (图~\ref{cd:five}):

\begin{figure}[htbp]
    \centering
    \begin{tikzcd}
        A_1 \arrow[r] \arrow[d, "\alpha_1"] & A_2 \arrow[r] \arrow[d, "\alpha_2"] & A_3 \arrow[r] \arrow[d, "\alpha_3"] & A_4 \arrow[r] \arrow[d, "\alpha_4"] & A_5 \arrow[d, "\alpha_5"] \\
        B_1 \arrow[r] & B_2 \arrow[r] & B_3 \arrow[r] & B_4 \arrow[r] & B_5
    \end{tikzcd}
    \caption{五引理}
    \label{cd:five}
\end{figure}

蛇引理 (图~\ref{cd:snake}):

\begin{figure}[htbp]
    \centering
    \begin{tikzcd}
         & 0 \arrow[d] & 0 \arrow[d] & 0 \arrow[d] & \\
        0 \arrow[r] & A' \arrow[r] \arrow[d, "f'"] & A \arrow[r] \arrow[d, "f"] & A'' \arrow[r] \arrow[d, "f''"] & 0 \\
        0 \arrow[r] & B' \arrow[r] \arrow[d, "g'"] & B \arrow[r] \arrow[d, "g"] & B'' \arrow[r] \arrow[d, "g''"] & 0 \\
         & C' \arrow[r] \arrow[d] & C \arrow[r] \arrow[d] & C'' \arrow[r] \arrow[d] & 0 \\
         & 0 & 0 & 0 &
    \end{tikzcd}
    \caption{蛇引理}
    \label{cd:snake}
\end{figure}

同调长正合列 (图~\ref{cd:les}):

\begin{figure}[htbp]
    \centering
    \begin{tikzcd}
        \cdots \arrow[r] & H_n(A) \arrow[r, "i_*"] & H_n(B) \arrow[r, "p_*"] & H_n(C) \arrow[r, "\partial"] & H_{n-1}(A) \arrow[r] & \cdots
    \end{tikzcd}
    \caption{同调长正合列}
    \label{cd:les}
\end{figure}

交换图引用测试: \autoref{cd:tensor}, \autoref{cd:free}, \autoref{cd:nat}, \autoref{cd:five}, \autoref{cd:snake}, \autoref{cd:les}.
